\SourceLeaf{344}

\Needspace{7\baselineskip}\begin{center}
[II.]

\textbf{[CIRCLE THROUGH FOUR POINTS.]}

[1.]
\end{center}

The circle through the points corresponding to the complex quantities $a$, $b$, $c$ with conjugates $A$, $B$, $C$ has as its centre the point corresponding to

\[
\begin{aligned}
&\frac{Aa(b-c)+Bb(c-a)+Cc(a-b)}{A(b-c)+B(c-a)+C(a-b)} \\
&= \frac{(Aa-Bb)(a-c)-(a-b)(Aa-Cc)}{(A-B)(a-c)-(a-b)(A-C)};
\end{aligned}
\]

more simply, when $a=A=0$:

\[
\frac{(B-C)bc}{Bc-Cb}.
\]

\Needspace{7\baselineskip}\begin{center}
[2.]
\end{center}

The condition for four points $0$, $a$, $b$, $c$ with conjugates
$0$, $A$, $B$, $C$ to lie on one and the same circle is the equation

\[
bcAB+caBC+abCA-abBC-bcCA-caAB = 0
\]
or more simply

\[
\frac{1}{aC}+\frac{1}{bA}+\frac{1}{cB}-\frac{1}{cA}-\frac{1}{aB}-\frac{1}{bC} = 0
\]
or also

\[
c(a-b)B(A-C) = b(a-c)C(A-B)
\]
etc.

\Needspace{7\baselineskip}\begin{center}
[3.]
\end{center}

The points reciprocal to three of those points with respect to the fourth
lie on one and the same straight line.

This theorem is included in the more general one:

\begin{quote}
“All points reciprocal to the points of a circle with respect to a given
point lie on one and the same circle”.
\end{quote}
